
\begin{equation}\tag{4.5}
 c_i(\underline{O}_Y)=(-1)^{i-1}(i-1)!\,\cl(Y)\quad \bmod
\text{ torsion},
\end{equation}
où dans le premier membre $\underline{O}_Y$ est considéré comme un faisceau cohérent sur $X$, définissant donc un élément de $K^\bullet(X)$, et dans le deuxième membre $\mathrm{c}\ell(Y)$ désigne la classe de $Y$ dans $A^i(X)$. Comme $A^i(X)$ ne change pas en enlevant de $X$ une partie fermée de codimension $>i$, on est ramené au cas où $Y$ est régulier, donc lisse sur le corps de base $k$ si ce dernier est parfait. Mais dans ce dernier cas, la formule (4.5) n'est autre que la formule de Riemann--Roch pour l'inclusion $Y\to X$ et l'élément $1$ de $K^\bullet(Y)$, formule écrite en utilisant la théorie des classes de Chern à valeurs dans l'anneau de Chow $A(X)$. Cette formule est établie dans [2] sous les conditions envisagées ici, par une méthode très voisine de celle utilisée dans ce séminaire pour la formule analogue utilisant $\operatorname{Gr}^\bullet(X)$. Cela établit donc, du moins pour $k$ parfait, la formule (4.4) et par suite le fait que $\phi$ soit un isomorphisme mod. torsion (qui entraîne à son tour, à posteriori, que la formule de Riemann--Roch ``avec dénominateurs'', pour un morphisme propre de schémas algébriques lisses sur $k$, est essentiellement la même, qu'on l'écrive en utilisant $A(X)$ comme dans [2], ou en utilisant $\operatorname{Gr}^\bullet(X)$ comme dans ce séminaire !). Le cas où on ne fait pas d'hypothèse sur $k$ se ramène aisément au cas où $k$ est algébriquement clos, en utilisant le fait que pour une extension finie $k'$ de degré $n$ de $k$, le noyau de l'application naturelle $A(X)\xrightarrow{f^*}A(X\otimes_k k')$ est annulé par $n$ (comme on voit en utilisant l'homomorphisme trace $f_*:A(X\otimes_k k')\to A(X)$, satisfaisant $f_*f^*=n\,\id$).\SGAAuthorFootText{*}{Elle est donc établie grâce à P. JOUANOLOU pour $X$ lisse quasi-projectif sur un corps (cf. note p.5).}

\textbf{4.3.} Il semble assez plausible que les formules (4.4) soient vraies telles quelles et non seulement modulo torsion, ou ce qui revient au même, que la formule 4.5. soit vraie non seulement modulo torsion. Cette question est un cas très particulier de celle de la validité dans $A(X)$ de la ``formule de Riemann--Roch sans dénominateurs'' (3.1), pour une immersion $Y\to X$, (savoir le cas où $i=d$). On peut considérer la validité de ce cas particulier comme assez plausible. En effet, nous avons signalé déjà au n\textsuperscript{o} 3 que lorsque $k$ est de caractéristique nulle, la formule de Riemann--Roch sans dénominateurs est valable tout au moins dans $\Gr^{\bullet}_{\mathrm{top}}(X)$, ce qui implique dans ce cas que les deux membres de (4.5) ont même image dans $\Gr^{\bullet}_{\mathrm{top}}(X)$, ou ce qui revient au même, que la deuxième des formules (4.4) est valable; ce résultat s'étendra de lui-même au cas de la caractéristique quelconque, une fois 
